%% kapa: exam=1 round=34 year=2023 period=1 form=A subject=economics subject_ko=경제학원론 seq=25 number=65
%% source: Q-Net 감정평가사 기출문제 (cst00309) artlSeq=5211789 file=raw/1008/5211789/2023년도 제34회 감정평가사 1차시험 1교시.pdf
%% conversion: kordoc 4.24.51 -> tools/gen1.py (auto); review=proofread-sonnet; flags=none
\begin{question}{65}
소비와 투자에 관한 설명으로 옳지 않은 것은?
\begin{choices}
  \item 소비수요는 사전적으로 계획된 소비를 말한다.
  \item 고전학파는 투자가 이자율이 아니라 소득에 의해 결정된다고 주장한다.
  \item 케인즈(J. Keynes)에 의하면 소비수요를 결정하는 중요한 요인은 현재의 절대소득이다.
  \item 독립투자수요는 내생변수와 관계없이 외생적으로 결정된다.
  \item 평균소비성향은 소비를 소득으로 나누어 계산한다.
\end{choices}
\end{question}
\answer{2}
